\subsection{Результаты в среде Ant}

Экспериментальное исследование в среде симуляции четырехногого робота (Ant-v4) стало завершающим и наиболее технически сложным этапом практической части работы. Данная среда, базирующаяся на физическом движке MuJoCo, отличается высокоразмерным непрерывным пространством состояний и действий, что предъявляет экстремальные требования к архитектуре когнитивного агента.


\subsubsection{Обоснование выбора модальности обучения}

В отличие от сред CartPole и LunarLander, где проводился сравнительный анализ двух алгоритмов, моделирование в среде Ant-v4 осуществлялось исключительно с использованием нейроэволюционного подхода (NEAT). Данное ограничение было введено намеренно и обусловлено фундаментальными математическими ограничениями классического алгоритма глубокого Q-обучения (DQN).

Поскольку управление «муравьем» требует подачи непрерывных сигналов на 8 независимых суставов, применение DQN потребовало бы искусственной дискретизации пространства действий. Даже при минимальном разбиении (например, на 3 состояния: тянуть назад, покой, тянуть вперед) алгоритм столкнулся бы с проблемой комбинаторного взрыва — $3^8 = 6561$ уникальная комбинация действий на каждом такте. Нейроэволюционный алгоритм NEAT, напротив, способен напрямую генерировать векторы непрерывных управляющих сигналов, что делает его оптимальным выбором для задач робототехники в рамках данного исследования.


\subsubsection{Эволюция функции вознаграждения (Reward Shaping) и адаптация среды}

Процесс обучения агента в среде Ant-v4 потребовал глубокого инженерного вмешательства в механику симуляции. В ходе первичных запусков с базовыми настройками среды был зафиксирован феномен стагнации приспособленности (отрицательные значения от -30 до 0). Анализ логов показал, что агент попадал в локальный оптимум, который можно охарактеризовать как «отказ от действий».

Для достижения итогового результата — биомеханически корректного и эффективного движения — был применен метод последовательного усложнения функции вознаграждения (Reward Shaping):

\paragraph{Решение проблемы «штрафа за энергию» (ctrl\_cost):} В базовой версии среды любые резкие сигналы на суставы карались жестким штрафом, который превышал награду за выживание. Эволюция NEAT реагировала на это радикально — алгоритм намеренно «ломал» нейросеть (уменьшая её размерность), чтобы выдавать нулевые сигналы и не получать штрафы. Для преодоления этой проблемы весовые коэффициенты штрафов в коде инициализации среды были снижены, что позволило агенту начать первичные исследования моторики без риска мгновенного завершения эпизода.

\paragraph{Внедрение «умного таймера» стабилизации:} Муравей появлялся в воздухе и падал на поверхность. Базовый код убивал агента за отсутствие скорости (таймер лени) еще до того, как тот успевал сгруппироваться после падения. Была разработана система «умного таймера», предоставляющая агенту запас времени (150 тактов) на стабилизацию после приземления и поощряющая любое, даже самое незначительное, продвижение вперед (установка собственного рекорда дистанции).

\paragraph{Получение биомеханически правильной походки:} После того как агент научился не погибать на старте и хаотично перебирать конечностями (простая функция награды), система вознаграждения была переориентирована на максимизацию проекции скорости вектора центра масс по оси X при минимизации энергозатрат.




\subsubsection{Практические выводы}

Благодаря поэтапной модификации функции вознаграждения, алгоритм NEAT успешно синтезировал топологию нейронной сети, способную управлять сложной многозвенной системой.

Эволюционный процесс привел к отказу от хаотичных, судорожных движений в пользу скоординированной походки. Агент научился использовать инерцию собственного тела, синхронизируя работу 8 суставов для поддержания баланса и минимизации контакта корпуса с землей. Итоговая модель нейросети оказалась поразительно компактной (учитывая сложность задачи), продемонстрировав способность нейроэволюции находить элегантные решения в многомерных пространствах без избыточного наращивания вычислительных мощностей.